
This work investigates whether LoRA adapter geometry encodes task similarity by analyzing the column spaces of learned B matrices under controlled experimental conditions. A total of 40 adapters were trained on 8 tasks spanning reasoning and natural language understanding categories, using identical base model (TinyLlama-1.1B-Chat-v1.0), hyperparameters, and training procedures. The primary finding is that adapters trained on semantically similar tasks exhibit smaller pairwise Grassmann distances than adapters trained on dissimilar tasks (Cohen's d = 0.765, p = 8.63 $\times 10^{-28})$. This geometric similarity correlates with FLAN task taxonomy (Spearman $\rho$ = 0.389, p = 1.29 $\times 10^{-29})$. Contrary to initial expectations, the clustering signal is distributed uniformly across all transformer layer types (attention and MLP) rather than concentrated in specific layers. A control analysis reveals that within-task variance across random seeds is substantial (ratio = 0.89 relative to within-category variance), indicating that training stochasticity contributes meaningfully to adapter geometry. These findings establish that task-category clustering exists in LoRA adapter weight space under controlled conditions, though practical utility for fine-grained task discrimination may be limited by training variance.
